\section{After Lefebvre: Continuity without Settlement}

Marcel Lefebvre died on 25 March 1991. The Society survived the founder because he had already transferred daily government to superiors and sacramental continuity to four bishops. Franz Schmidberger served as superior general until 1994; Bernard Fellay then held the office for twenty-four years. The bishops did not receive territorial dioceses or a claim of ordinary jurisdiction in the Society's own structure. They functioned as auxiliaries for ordination and confirmation while the superior general governed the priestly body. That arrangement preserved the claim that the 1988 act was provision for emergency, not creation of a parallel territorial hierarchy, while leaving the basic question of mission unresolved.

The Pontifical Commission \emph{Ecclesia Dei} developed the other institutional legacy of 1988. It assisted clergy, religious, and faithful attached to older liturgical forms who sought recognized structures. Its work did not make the SSPX itself canonical; it created alternatives and maintained channels through which reconciliation could be explored.

\subsection{Formal adherence and ordinary attendance}

A 24 August 1996 note of the Pontifical Council for Legislative Texts addressed the schismatic situation associated with Lefebvre and the meaning of formal adherence to it. Until a material change could be established, the note said, the movement was to be considered schismatic; ministerial activity within it was strong evidence of formal adherence. For an individual lay person, however, the note required internal and external elements and evaluation of intention and conduct. It explicitly prevented a mechanical equation between occasional participation in SSPX liturgy and formal adherence, while treating sustained conduct that embodied rejection of the Roman Pontiff's authority differently from mere attendance.

This distinction would become newly important in 2026. It is also a model for historical precision: penalties attached to named consecrating and consecrated bishops in 1988, the canonical condition of SSPX clergy, and the possible culpability of a lay participant were related but not identical questions.

\subsection{The Jubilee pilgrimage of 2000}

During the Great Jubilee, SSPX bishops, priests, and faithful made a pilgrimage to Rome. The gesture did not constitute reconciliation, yet it made visible a paradox at the center of the Society: it disputed Roman acts while making attachment to Rome part of its identity. Contacts intensified after the pilgrimage. Both parties could interpret the same act favorably for different reasons---the Holy See as a possible opening to communion, the SSPX as a profession of Roman Catholic identity without surrender of its objections.

\begin{studybox}{Institutional continuity after 1991}
The death of a charismatic founder often tests whether a movement is an organization or an episode. The SSPX retained a general government, houses, seminaries, recurrent ordinations, four bishops, and an intergenerational constituency. Its survival established sociological durability. It did not by itself answer the canonical question that the founder's acts had intensified.
\end{studybox}

The longer the interim lasted, the more difficult a settlement became. Priests formed entirely within SSPX institutions inherited not only an older liturgy but a developed memory of suppression, penalties, negotiations, and providential survival. Roman concessions could be welcomed as admissions that the Society's earlier diagnosis had merit, while Roman authorities could intend them as pastoral steps toward a unity still lacking. The same measure could therefore reduce immediate hardship without resolving the interpretive conflict.
